\section{Conclusions}
Electrical measurements reconstruct terminal-current flux samples, while
co-energy models a conservative magnetic storage subsystem in its conjugate
coordinates. Interpreting those samples as its gradients requires additional
coordinate, state, and admissibility evidence. Path independence explains
integrability and gauge freedom; direct gradient fitting uses noisy samples
jointly without first integrating them along an arbitrary path.

An even co-energy in q current yields d-even and q-odd flux. Its Hessian is
the differential-inductance matrix, with reciprocal cross terms and
operating-point-dependent curvature that accommodates saturation and
cross-saturation. A symmetric scalar B-spline potential retains these
relationships by construction; independent fits do not enforce them unless
explicit constraints couple the two components.

The synthetic illustration separates flux accuracy, derivative quality, and
physical consistency. It shows that parity constraints alone can leave curl,
that a common potential eliminates curl, and that symmetry and integrability
still do not guarantee positive curvature or truth. Regularization, data
support, robust fitting, and additional physical constraints remain necessary
engineering choices. Once the admissible magnetic structure is defined,
departures can be preserved as diagnostic information rather than absorbed
blindly by an interpolator. The companion paper studies that next question.

The real-data stress test shows why these distinctions matter. A resistance
hypothesis improves selected high-current parity but worsens all-pair d/vector
residuals and relative path disagreement. The symmetry and path equivalents
disagree even with exact P1 integration; their gap is not a trapez-quadrature
effect. Independent P1 flux interpolation can itself break integrability,
and local derivatives depend strongly on mesh geometry. Terminal and
magnetizing currents, normalization, and unmodeled loss or measurement states
remain open. A later potential fit would be a projection onto a declared
model class whose residual must remain visible. These data identify neither
a true winding resistance nor an iron-loss explanation or successful real-data
co-energy reconstruction.
